% Section 2.9, from lectures/L02/appendix/b_exercises.md.

\section{Uppgifter}
\label{c2:sec:uppgifter}

\begin{aside}
\textbf{Så arbetar du med uppgifterna.} Del~1 och del~2 räknas under lektionen: först på egen
hand, några minuter per uppgift, sedan gemensamt. Del~3 är extrauppgifter att träna vidare på
efter lektionen. De flesta går att räkna i huvudet eller med papper och penna.

\facitnotis{L02}
\end{aside}

% ------------------------------------------------------------------------------------------
\uppgiftsdel{Del 1}{Repetitionsuppgifter}

\uppgift{1.1}{Parallellkoppling av tre motstånd}
Tre motstånd $R_1 = 6\,\Omega$, $R_2 = 3\,\Omega$ och $R_3 = 2\,\Omega$ är parallellkopplade.

\bild[0.52\linewidth]{lectures/L02/appendix/images/1.1_circuit.png}

Beräkna den totala resistansen $R_{\text{TOT}}$, som ges av
\[
  \frac{1}{R_{\text{TOT}}} = \frac{1}{R_1} + \frac{1}{R_2} + \frac{1}{R_3}.
\]

\uppgift{1.2}{Verkningsgrad och effekt}
En transformator har verkningsgraden $\eta = 92\,\%$ och ineffekten $P_{\text{in}} = 500\,\text{W}$.

\bild[0.4\linewidth]{lectures/L02/appendix/images/1.2_circuit.png}

\begin{parts}
  \item Beräkna uteffekten $P_{\text{ut}}$.
  \item Hur stor effekt $P_{\text{förl}}$ förloras som värme?
\end{parts}

\needspace{18\baselineskip}
\uppgift{1.3}{Räkneordning i kretsformel}
Beräkna spänningsdelarens utspänning $U_{\text{ut}}$ med formeln nedan då
$U_{\text{in}} = 12\,\text{V}$, $R_1 = 8\,\Omega$ och $R_2 = 4\,\Omega$.

\bild[0.44\linewidth]{lectures/L02/appendix/images/1.3_circuit.png}

\[
  U_{\text{ut}} = U_{\text{in}} \times \frac{R_2}{R_1 + R_2}
\]

% ------------------------------------------------------------------------------------------
\uppgiftsdel{Del 2}{Nytt stoff}

\uppgift{2.1}{Förenkling av algebraiska uttryck}
Förenkla följande uttryck.
\begin{delar}(2)
  \task $5I_1 + 3I_2 - 2I_1 + 7I_2$
  \task $4R_1 - 3R_2 + R_1 + 5R_2 - 2R_1$
  \task $3I^2 - 2I + 4I^2 + 5I - 1$
  \startnewitemline
  \task* Sätt in $I = 2$ i både det ursprungliga och det förenklade uttrycket i c) och kontrollera
    att de ger samma värde.
\end{delar}

\needspace{24\baselineskip}
\uppgift{2.2}{Distributivlagen i kretsar}
Motståndet $R$ i kretsen nedan genomflyts av summan av de tre grenströmmarna $I_1$, $I_2$ och
$I_3$.

\bild[0.7\linewidth]{lectures/L02/appendix/images/2.2_circuit.png}

Spänningsfallet över $R$ ges därmed av
\[
  U_R = RI = R(I_1 + I_2 + I_3).
\]
\begin{parts}
  \item Multiplicera ut parentesen och ange varje term med enhet.
  \item Beräkna $U_R$ med den ursprungliga formen då $R = 100\,\Omega$, $I_1 = 20\,\text{mA}$,
    $I_2 = 30\,\text{mA}$ och $I_3 = 50\,\text{mA}$.
  \item Beräkna $U_R$ en gång till med det utmultiplicerade uttrycket och kontrollera att svaren
    stämmer överens.
\end{parts}

\uppgift{2.3}{Faktorisering}
Faktorisera följande uttryck.
\begin{delar}(4)
  \task $6I + 9RI$
  \task $P_1^2 - P_2^2$
  \task $U^2 - 10U + 25$
  \task $4R^2 - 1$
\end{delar}

\uppgift{2.4}{Gemensam faktor: total effekt i en seriekrets}
Tre motstånd $R_1 = 120\,\Omega$, $R_2 = 180\,\Omega$ och $R_3 = 200\,\Omega$ är seriekopplade och
genomflyts därmed av samma ström $I = 50\,\text{mA}$.

\bild[0.52\linewidth]{lectures/L02/appendix/images/2.4_circuit.png}

Effekten i respektive motstånd ges av $P_k = R_k I^2$.
\begin{parts}
  \item Skriv ett uttryck för den totala effekten $P_{\text{TOT}} = P_1 + P_2 + P_3$ och bryt ut den
    gemensamma faktorn.
  \item Beräkna $P_{\text{TOT}}$ med det faktoriserade uttrycket.
  \item Beräkna $P_1$, $P_2$ och $P_3$ var för sig och kontrollera att summan blir densamma. Vilken
    väg krävde minst räknearbete?
\end{parts}

\uppgift{2.5}{Konjugatregeln: skillnad i effekt}
Ett värmeelement med resistansen $R$ matas först med spänningen $U_1$ och sedan med $U_2$.

\bild[0.35\linewidth]{lectures/L02/appendix/images/2.5_circuit.png}

Effekten ges av $P = \dfrac{U^2}{R}$.
\begin{parts}
  \item Visa att effektskillnaden kan skrivas
    \[
      P_1 - P_2 = \frac{(U_1 + U_2)(U_1 - U_2)}{R}.
    \]
  \item Beräkna $P_1 - P_2$ \textbf{utan} miniräknare då $U_1 = 24\,\text{V}$,
    $U_2 = 20\,\text{V}$ och $R = 8\,\Omega$.
  \item Kontrollera svaret genom att beräkna $P_1$ och $P_2$ var för sig.
\end{parts}

\uppgift{2.6}{Kvadreringsregeln: effekt vid en strömändring}
Genom ett motstånd $R = 10\,\Omega$ flyter strömmen $I_0 = 2\,\text{A}$. Strömmen ökar med
$\Delta I = 0{,}1\,\text{A}$, så att den nya strömmen blir $I_0 + \Delta I$.

\bild[0.35\linewidth]{lectures/L02/appendix/images/2.6_circuit.png}

Effekten ges av $P = RI^2$.
\begin{parts}
  \item Multiplicera ut $P = R(I_0 + \Delta I)^2$ med kvadreringsregeln.
  \item Beräkna den nya effekten med det utmultiplicerade uttrycket.
  \item Hur stor är effektökningen $\Delta P$ jämfört med $P_0 = RI_0^2$?
  \item Hur stort blir felet om man struntar i termen $R\,\Delta I^2$? Ange felet i procent av
    $\Delta P$.
\end{parts}

\needspace{22\baselineskip}
\uppgift{2.7}{Effektuttryck och Ohms lag}
Effekten som utvecklas i motståndet nedan kan skrivas $P = RI^2$.

\bild[0.31\linewidth]{lectures/L02/appendix/images/2.7_circuit.png}

\begin{parts}
  \item Visa att $P = RI^2$ kan faktoriseras till $P = (RI) \cdot I$ och identifiera vilken storhet
    $RI$ representerar.
  \item Sätt in $I = \dfrac{U}{R}$ i $P = RI^2$ och visa att uttrycket kan förenklas till
    $P = \dfrac{U^2}{R}$.
  \item Ett motstånd $R = 6\,\Omega$ matas med $U = 12\,\text{V}$. Beräkna effekten med alla tre
    formerna $P = UI$, $P = RI^2$ och $P = \dfrac{U^2}{R}$ och kontrollera att de ger samma svar.
\end{parts}

% ------------------------------------------------------------------------------------------
\uppgiftsdel{Del 3}{Extrauppgifter}
Uppgifterna nedan är extra träning och görs med fördel på egen hand efter lektionen.

\uppgift{3.1}{Termer, faktorer och koefficienter}
Betrakta uttrycket $7R^2 - 4R + 9$.
\begin{parts}
  \item Hur många termer består uttrycket av?
  \item Ange koefficienten till $R^2$ respektive till $R$.
  \item Ange konstanttermen.
  \item Vilka faktorer ingår i termen $7R^2$?
  \item Beräkna uttryckets värde för $R = 2$.
\end{parts}

\uppgift{3.2}{Förenkling med parenteser}
Förenkla.
\begin{delar}(2)
  \task $2(3I + 4) - 3(I - 2)$
  \task $-(U - 5) + 2(U + 1)$
  \task $5R_1 - \left[2R_1 - (R_1 + 3)\right]$
  \task $3(2P - 1) - 2(3P - 4)$
\end{delar}

\uppgift{3.3}{Multiplicera ut}
Multiplicera ut och förenkla.
\begin{delar}(3)
  \task $(R + 3)(R + 5)$
  \task $(2I - 1)(I + 4)$
  \task $(U + 2)(U - 2)$
  \task $(3R - 2)(3R - 2)$
  \task $2I(I + 3) - I(2I - 1)$
\end{delar}

\uppgift{3.4}{Kvadreringsreglerna}
Utveckla.
\begin{delar}(4)
  \task $(I + 4)^2$
  \task $(R - 6)^2$
  \task $(2U + 3)^2$
  \task $(5 - I)^2$
  \task* Visa med $a = 3$ och $b = 4$ att $(a + b)^2 \neq a^2 + b^2$.
\end{delar}

\uppgift{3.5}{Konjugatregeln som huvudräkningstrick}
Använd $(a + b)(a - b) = a^2 - b^2$ och räkna \textbf{utan} miniräknare.
\begin{delar}(4)
  \task $102 \times 98$
  \task $45 \times 35$
  \task $2{,}1 \times 1{,}9$
  \task $51^2 - 49^2$
  \task* Spänningen över ett motstånd $R = 5\,\Omega$ ändras från $U_1 = 25\,\text{V}$ till
    $U_2 = 15\,\text{V}$. Beräkna $P_1 - P_2 = \dfrac{U_1^2 - U_2^2}{R}$ utan att kvadrera något
    tal.
\end{delar}

\uppgift{3.6}{Faktorisera}
Faktorisera så långt som möjligt.
\begin{delar}(3)
  \task $8R + 12$
  \task $RI^2 + RI$
  \task $U^2 - 49$
  \task $I^2 + 12I + 36$
  \task $9R^2 - 16$
  \task $2U^2 - 8$
\end{delar}

\uppgift{3.7}{Förkorta algebraiska bråk}
Faktorisera och förkorta. Antag att nämnaren inte är noll.
\begin{delar}(3)
  \task $\dfrac{6RI}{3R}$
  \task $\dfrac{R^2 - 9}{R + 3}$
  \task $\dfrac{2U^2 + 4U}{2U}$
  \task $\dfrac{I^2 - 10I + 25}{I - 5}$
  \task $\dfrac{4R^2 - 1}{2R - 1}$
\end{delar}

\needspace{9\baselineskip}
\uppgift{3.8}{Insättning i effektformlerna}
Ett motstånd $R = 8\,\Omega$ genomflyts av strömmen $I = 1{,}5\,\text{A}$.
\begin{parts}
  \item Beräkna spänningen $U = RI$.
  \item Beräkna effekten med $P = RI^2$.
  \item Beräkna effekten med $P = \dfrac{U^2}{R}$ och kontrollera att svaret stämmer.
  \item Strömmen fördubblas. Med vilken faktor ändras effekten? Motivera algebraiskt.
\end{parts}

\needspace{13\baselineskip}
\uppgift{3.9}{Spänningsdelarens algebra}
För två seriekopplade motstånd gäller
\[
  U_1 = U_{\text{in}} \times \frac{R_1}{R_1 + R_2}
  \qquad \text{och} \qquad
  U_2 = U_{\text{in}} \times \frac{R_2}{R_1 + R_2}.
\]
\begin{parts}
  \item Visa algebraiskt att $U_1 + U_2 = U_{\text{in}}$.
  \item Visa att $\dfrac{U_1}{U_2} = \dfrac{R_1}{R_2}$.
  \item Vad blir $U_1$ om $R_1 = R_2$?
  \item Beräkna $U_1$ och $U_2$ då $U_{\text{in}} = 15\,\text{V}$, $R_1 = 1\,\text{k}\Omega$ och
    $R_2 = 4\,\text{k}\Omega$.
\end{parts}

\uppgift{3.10}{Hitta felet}
Varje likhet nedan innehåller ett vanligt algebrafel. Förklara vad som är fel och skriv det
korrekta högerledet.
\begin{delar}(3)
  \task $(R + 2)^2 = R^2 + 4$
  \task $3(I - 2) = 3I - 2$
  \task $\dfrac{R + 4}{4} = R$
  \task $-(U - 3) = -U - 3$
  \task $2R \times 3R = 6R$
  \task $(2I)^2 = 2I^2$
\end{delar}
